% GENERATED FILE. Edit canonical lessons, then run npm run generate:books.
% canonical-source-hash: fnv1a64:80450bf1381e6421
% canonical-lessons: JA-C281-uwasa, JA-C281-tsutawaru, JA-C281-hiromaru, JA-C281-wadai, JA-C281-ousawagi

\chapter{Word Gets Around}
\label{ch:ja-word-gets-around}
\hlchaptermodality{281}

\begin{chapteropening}
\textbf{By the end of this chapter:} \emph{I can say how news and rumours spread.}

Complete the last lesson of chapter 281: I can say how news and rumours spread.
\end{chapteropening}

\section[\texorpdfstring{\ja{うわさ} (uwasa)}{うわさ (uwasa)}]{\ja{うわさ} (uwasa) — a rumour; gossip}

\label{lesson:JA-C281-uwasa}



\begin{quote}
Before the new one: say the Japanese for a reporter, then the Japanese for to hand out.
\end{quote}

\subsection*{You'll want to know: \ja{うわさ}}
\textbf{\ja{うわさ}} — \emph{uwasa} — \textquotedblleft{}a rumour; gossip\textquotedblright{}.

\begin{quote}
\textbf{\ja{その} \ja{うわさは} \ja{ほんとうですか。}} — \emph{sono uwasa wa hontō desu ka} — \textquotedblleft{}Is the rumour true?\textquotedblright{}
\end{quote}

\begin{cousinweb}
\textbf{\ja{うわさ}} (\emph{uwasa}) is a native Japanese word (wago), inherited rather than borrowed from Chinese; its kanji was attached later as a kun'yomi (JMdict).
\end{cousinweb}

\begin{grammarlens}[title={What you've built}]
One more word for telling what happened.
\end{grammarlens}

\subsection*{Your turn}
\begin{itemize}
\raggedright
  \item \emph{Say it:} \emph{uwasa}
  \item \emph{Say it:} \emph{uwasa}, once more
  \item \emph{Say it:} \emph{sono uwasa wa hontō desu ka} — Is the rumour true?
  \item \emph{Recall:} say the Japanese for a reporter, then the Japanese for to hand out, then say \emph{uwasa} again
\end{itemize}

\begin{tcolorbox}[breakable,colback=teal!4,colframe=teal!35!black,title={Before you move on}]
What does \ja{うわさ} mean? (\textquotedblleft{}A rumour; gossip\textquotedblright{}.) Say it once more.
\end{tcolorbox}
\section[\texorpdfstring{\ja{つたわる} (tsutawaru)}{つたわる (tsutawaru)}]{\ja{つたわる} (tsutawaru) — to be passed on; to spread (of news)}

\label{lesson:JA-C281-tsutawaru}



\begin{quote}
Before the new one: say the Japanese for to hand out, then the Japanese for a rumour.
\end{quote}

\subsection*{You'll want to know: \ja{つたわる}}
\textbf{\ja{つたわる}} — \emph{tsutawaru} — \textquotedblleft{}to be passed on; to spread (of news)\textquotedblright{}.

\begin{quote}
\textbf{\ja{じけんは} \ja{すぐ} \ja{まちの} \ja{みんなに} \ja{つたわりました。}} — \emph{jiken wa sugu machi no minna ni tsutawarimashita} — \textquotedblleft{}News of the incident soon reached the whole town.\textquotedblright{}
\end{quote}

Polite \textbf{\ja{つたわります}}, \emph{tsutawarimasu}; a godan verb. Past (polite): \textbf{\ja{うわさが} \ja{つたわりました。}}, \emph{uwasa ga tsutawarimashita}.

\begin{cousinweb}
\textbf{\ja{つたわる}} (\emph{tsutawaru}) is a native Japanese word (wago), inherited rather than borrowed from Chinese; its kanji was attached later as a kun'yomi (JMdict). It is the intransitive partner of \emph{tsutaeru}, \textquotedblleft{}to pass on\textquotedblright{}.
\end{cousinweb}

\begin{grammarlens}[title={What you've built}]
One more word for telling what happened.
\end{grammarlens}

\subsection*{Your turn}
\begin{itemize}
\raggedright
  \item \emph{Say it:} \emph{tsutawaru}
  \item \emph{Say it:} \emph{tsutawaru}, once more
  \item \emph{Say it:} \emph{jiken wa sugu machi no minna ni tsutawarimashita} — News of the incident soon reached the whole town.
  \item \emph{Recall:} say the Japanese for to hand out, then the Japanese for a rumour, then say \emph{tsutawaru} again
\end{itemize}

\begin{tcolorbox}[breakable,colback=teal!4,colframe=teal!35!black,title={Before you move on}]
What does \ja{つたわる} mean? (\textquotedblleft{}To be passed on; to spread (of news)\textquotedblright{}.) Say it once more.
\end{tcolorbox}
\section[\texorpdfstring{\ja{ひろまる} (hiromaru)}{ひろまる (hiromaru)}]{\ja{ひろまる} (hiromaru) — to spread widely; to get about}

\label{lesson:JA-C281-hiromaru}



\begin{quote}
Before the new one: say the Japanese for a rumour, then the Japanese for to be passed on.
\end{quote}

\subsection*{You'll want to know: \ja{ひろまる}}
\textbf{\ja{ひろまる}} — \emph{hiromaru} — \textquotedblleft{}to spread widely; to get about\textquotedblright{}.

\begin{quote}
\textbf{\ja{その} \ja{うわさは} \ja{がっこうで} \ja{ひろまっています。}} — \emph{sono uwasa wa gakkō de hiromatte imasu} — \textquotedblleft{}The rumour is going round the school.\textquotedblright{}
\end{quote}

Polite \textbf{\ja{ひろまります}}, \emph{hiromarimasu}; a godan verb. Past (polite): \textbf{\ja{うわさが} \ja{ひろまりました。}}, \emph{uwasa ga hiromarimashita}.

\begin{cousinweb}
\textbf{\ja{ひろまる}} (\emph{hiromaru}) is a native Japanese word (wago), inherited rather than borrowed from Chinese; its kanji was attached later as a kun'yomi (JMdict).
\end{cousinweb}

\begin{grammarlens}[title={What you've built}]
One more word for telling what happened.
\end{grammarlens}

\subsection*{Your turn}
\begin{itemize}
\raggedright
  \item \emph{Say it:} \emph{hiromaru}
  \item \emph{Say it:} \emph{hiromaru}, once more
  \item \emph{Say it:} \emph{sono uwasa wa gakkō de hiromatte imasu} — The rumour is going round the school.
  \item \emph{Recall:} say the Japanese for a rumour, then the Japanese for to be passed on, then say \emph{hiromaru} again
\end{itemize}

\begin{tcolorbox}[breakable,colback=teal!4,colframe=teal!35!black,title={Before you move on}]
What does \ja{ひろまる} mean? (\textquotedblleft{}To spread widely; to get about\textquotedblright{}.) Say it once more.
\end{tcolorbox}
\section[\texorpdfstring{\ja{わだい} (wadai)}{わだい (wadai)}]{\ja{わだい} (wadai) — a topic of conversation}

\label{lesson:JA-C281-wadai}



\begin{quote}
Before the new one: say the Japanese for to be passed on, then the Japanese for to spread widely.
\end{quote}

\subsection*{You'll want to know: \ja{わだい}}
\textbf{\ja{わだい}} — \emph{wadai} — \textquotedblleft{}a topic of conversation\textquotedblright{}.

\begin{quote}
\textbf{\ja{その} \ja{じけんが} \ja{わだいに} \ja{なっています。}} — \emph{sono jiken ga wadai ni natte imasu} — \textquotedblleft{}Everyone is talking about the incident.\textquotedblright{}
\end{quote}

\begin{cousinweb}
\textbf{\ja{わだい}} (\emph{wadai}) is Sino-Japanese (kango): two Chinese characters read with their on'yomi, readings borrowed from Chinese. They are \emph{wa}, \textquotedblleft{}talk\textquotedblright{} (Middle Chinese \emph{hwaejH}) and \emph{dai}, \textquotedblleft{}topic\textquotedblright{} (\emph{dej}) (KANJIDIC2; Middle Chinese in Baxter's transcription).
\end{cousinweb}

\begin{grammarlens}[title={What you've built}]
One more word for telling what happened.
\end{grammarlens}

\subsection*{Your turn}
\begin{itemize}
\raggedright
  \item \emph{Say it:} \emph{wadai}
  \item \emph{Say it:} \emph{wadai}, once more
  \item \emph{Say it:} \emph{sono jiken ga wadai ni natte imasu} — Everyone is talking about the incident.
  \item \emph{Recall:} say the Japanese for to be passed on, then the Japanese for to spread widely, then say \emph{wadai} again
\end{itemize}

\begin{tcolorbox}[breakable,colback=teal!4,colframe=teal!35!black,title={Before you move on}]
What does \ja{わだい} mean? (\textquotedblleft{}A topic of conversation\textquotedblright{}.) Say it once more.
\end{tcolorbox}
\section[\texorpdfstring{\ja{おおさわぎ} (ōsawagi)}{おおさわぎ (ōsawagi)}]{\ja{おおさわぎ} (ōsawagi) — an uproar; a commotion}

\label{lesson:JA-C281-ousawagi}



\begin{quote}
Before the new one: say the Japanese for to spread widely, then the Japanese for a topic of conversation.
\end{quote}

\subsection*{You'll want to know: \ja{おおさわぎ}}
\textbf{\ja{おおさわぎ}} — \emph{ōsawagi} — \textquotedblleft{}an uproar; a commotion\textquotedblright{}.

\begin{quote}
\textbf{\ja{えきで} \ja{おおさわぎに} \ja{なりました。}} — \emph{eki de ōsawagi ni narimashita} — \textquotedblleft{}There was a great commotion at the station.\textquotedblright{}
\end{quote}

\begin{cousinweb}
\textbf{\ja{おおさわぎ}} (\emph{ōsawagi}) is a native Japanese word (wago), inherited rather than borrowed from Chinese, built from \emph{ō}, \textquotedblleft{}big\textquotedblright{}, and \emph{sawagi}, \textquotedblleft{}noise\textquotedblright{} (from \emph{sawagu}) (JMdict, KANJIDIC2).
\end{cousinweb}

\begin{grammarlens}[title={What you've built}]
One more word for telling what happened.
\end{grammarlens}

\subsection*{Your turn}
\begin{itemize}
\raggedright
  \item \emph{Say it:} \emph{ōsawagi}
  \item \emph{Say it:} \emph{ōsawagi}, once more
  \item \emph{Say it:} \emph{eki de ōsawagi ni narimashita} — There was a great commotion at the station.
  \item \emph{Recall:} say the Japanese for to spread widely, then the Japanese for a topic of conversation, then say \emph{ōsawagi} again
\end{itemize}

\begin{tcolorbox}[breakable,colback=teal!4,colframe=teal!35!black,title={Before you move on}]
What does \ja{おおさわぎ} mean? (\textquotedblleft{}An uproar; a commotion\textquotedblright{}.) Say it once more.
\end{tcolorbox}
